\subsection{单轨道三场解及逆向验证}
重构公式将同一矩阵的无穷远系数与共振值分别转化为光场和介质. 以下结论同时给出三个物理场及其 Lax 相容性.

\begin{theorem}[单轨道三场解]\label{thm:direct-HMB}\label{prop:oneboundary}
设定理 \ref{thm:oneRHP} 的条件成立, $N$ 按 \eqref{eq:Nfactor} 定义, $n_{jk}=N(a)_{jk}$. 则物理场
\begin{equation}\label{eq:finalfields}
\begin{aligned}
\widetilde q&=-e^{-\ii\Omega t}(q_0+4bP_{12}),\\
\widetilde p&=\dfrac{e^{-\ii\Omega t}(N(a)R_\infty N(a)^\dagger)_{12}}{a^2+b^2},\qquad
\eta=\dfrac{(N(a)R_\infty N(a)^\dagger)_{11}}{a^2+b^2}
\end{aligned}
\end{equation}
在全体实 $x,t$ 上光滑, 属于定义 \ref{ass:fields} 中 $I=\mathbb R$ 的基本场类, 并满足原 HMB 系统 \eqref{eq:HMB}、Bloch 约束 \eqref{eq:C} 及两端同相边界 \eqref{eq:targetboundary}. 两个共振原像重构出同一个 Hermitian 迹零介质矩阵. 介质分量为
\begin{align}
\eta&=\dfrac{\eta_0}{a^2+b^2}\left[|n_{11}|^2-|n_{12}|^2
+\dfrac{\ii q_0}{a}(n_{11}\conj{n_{12}}-n_{12}\conj{n_{11}})\right],\label{eq:etacomponent}\\
e^{\ii\Omega t}\widetilde p&=\dfrac{\eta_0}{a^2+b^2}\left[n_{11}\conj{n_{21}}-n_{12}\conj{n_{22}}
+\dfrac{\ii q_0}{a}(n_{11}\conj{n_{22}}-n_{12}\conj{n_{21}})\right].\label{eq:pcomponent}
\end{align}
固定 $t$ 时, 三个场的端态误差为 $O(e^{-2\kappa|x|})$.
\end{theorem}
\begin{proof}
\par\medskip\noindent\textbf{三场重构.}\quad
在无穷远展开 \eqref{eq:explicitM}, 得
\[
M^{(1)}=H(2\ii bI-4\ii bP-\ii q_0\sigma_1-2\ii\kappa\sg)H^{-1}.
\]
于是规范光场 $q=\ii M^{(1)}_{12}=-q_0-4bP_{12}$. 在实共振点, 由 $Y\lambda\sg Y^{-1}=\ii U_\infty$ 及 $N(a)^{-1}=\dfrac{N(a)^\dagger}{a^2+b^2}$,
\begin{equation}\label{eq:direct-medium}
M(z_a)\dfrac{\eta_0\lambda(z_a)}{a}\sg M(z_a)^{-1}
=\dfrac{HN(a)R_\infty N(a)^\dagger H^{-1}}{a^2+b^2}=:R.
\end{equation}
右端与谱叶无关. $N(a)N(a)^\dagger=(a^2+b^2)I$, 因而该酉共轭保持 Hermitian 性、迹零及 $R_\infty$ 的特征值. 特别地,
$R^2=\eta_0^2\left(1+\dfrac{q_0^2}{a^2}\right)I$, 给出 Bloch 约束. 结合 \eqref{eq:original-fields} 得到 \eqref{eq:finalfields}; 矩阵相乘给出两条分量式.

\par\medskip\noindent\textbf{空间端态与基本场类.}\quad
当 $x\to\pm\infty$ 时, $|T|$ 分别趋于无穷和零, 因而
\begin{equation}\label{eq:Pends}
P_+=\dfrac{1}{1+h_0^2}\begin{pmatrix}1&-h_0\\-h_0&h_0^2\end{pmatrix},\qquad
P_-=\dfrac{1}{1+h_0^2}\begin{pmatrix}h_0^2&-h_0\\-h_0&1\end{pmatrix}.
\end{equation}
这给出 $q\to q_0$. 将端态代入 $N$, 有
\[
N_+Y=Y_{-q_0}\diag(\lambda-\ii\kappa,\lambda+\ii\kappa),\qquad
N_-Y=Y_{-q_0}\diag(\lambda+\ii\kappa,\lambda-\ii\kappa),
\]
其中 $Y_{-q_0}=I+\dfrac{\ii q_0}{z}\sigma_1$. 利用 $HY_{-q_0}H^{-1}=Y$, 得
\begin{equation}\label{eq:direct-M-ends}
M\longrightarrow Y\quad(x\to-\infty),\qquad
M\longrightarrow Y\diag(s_{22},s_{22}^{-1})\quad(x\to+\infty).
\end{equation}
对角因子与 $\sg$ 对易, 所以两端 $R\to R_\infty$. 由 $|T|=e^{2\kappa[x-\operatorname{Re}F(\ii b)t]}$, 端态误差及任意固定阶导数均按 $e^{-2\kappa|x|}$ 衰减, 对紧 $t$ 区间一致. $\Delta_->0$ 与 $a^2+b^2>0$ 保证全局光滑.


更具体地, 对任意紧区间 $J\subset\mathbb R$ 和任意固定非负整数 $r,s$,
对端态扰动求 $\partial_x^r\partial_t^s$ 后仍有
$C_{J,r,s}e^{-2\kappa|x|}$ 的统一界.
这是因为 $T_x=2\kappa T$, $T_t=-2\kappa F(\ii b)T$,
而各场是 $T,\conj T$ 的有理式, 分母由 \eqref{eq:delta-pair} 的正下界控制.
在两端分别用 $T$ 或 $T^{-1}$ 展开, 即得到同阶的导数估计.
该可积控制同时允许对加权范数使用支配收敛,
所以三场满足定义 \ref{ass:fields}, 其中 $I=\mathbb R$.


\par\medskip\noindent\textbf{矩阵对数导数与奇性消去.}\quad
对避开特殊点的 $z$ 定义
\begin{equation}\label{eq:inverse-log-lax}
U=(M_x-\ii\lambda M\sg)M^{-1},\qquad
V=(M_t+\ii\lambda F M\sg)M^{-1}.
\end{equation}
它们是 $\phi=Me^{\ii\theta\sg}$ 的两个对数导数. 换叶时 $\lambda$ 和 $\theta$ 变号, $\sigma_1$ 与 $\sg$ 反对易, 故两矩阵在 $z$ 和 $h(z)$ 取值相同, 是 $k$ 的单值有理矩阵.

在 $\zeta_j$ 附近, $\phi$ 的列留数为 $\Res\phi_2=c_j(0)\phi_1(\zeta_j)$. 常数 $c_j(0)$ 与 $x,t$ 无关, 故
\[
\phi=H_j(z)\left(I+\dfrac{c_j(0)e_1e_2^T}{z-\zeta_j}\right),\qquad
\det H_j(\zeta_j)=\gamma(\zeta_j)\ne0.
\]
右因子不依赖 $x,t$, 因而两个对数导数的离散极点可去; 共轭极点同理.

在分支点 $z_*=\pm\ii q_0$, 换叶约化使 $M(z_*)$ 消灭 $C(z_*)$ 的固定 $-1$ 特征空间. 该核不依赖 $x,t$, 因此 $M_x(z_*)$ 和 $M_t(z_*)$ 也消灭它. 又 $\lambda(z_*)=0$, $F$ 在此有限, \eqref{eq:inverse-log-lax} 的两个分子均消灭同一核. $\det M=\gamma$ 在 $z_*$ 有一阶零点, 而 $\operatorname{adj}M(z_*)$ 的像正是该核, 故两个商的分子抵消行列式的零点. 这证明分支点可去. 轮廓交点处 $\gamma\ne0$, 没有逆矩阵奇性. 在共振点, $U$ 解析, $V$ 的唯一留数为
\[
\Res_{k=a}V=\dfrac{\ii\eta_0\lambda(z_a)}{a}M(z_a)\sg M(z_a)^{-1}=\ii R.
\]

\par\medskip\noindent\textbf{恢复 Lax 系数并验证方程.}\quad
在第一叶无穷远, $M=I+O(k^{-1})$, $M_x,M_t=O(k^{-1})$. 由 \eqref{eq:reconq},
\[
U=-\ii k\sg+Q+O(k^{-1}).
\]
有限点均可去, Liouville 定理给出 $U=-\ii k\sg+Q$. 恒等式
\[
V=-FU+(M_t+FM_x)M^{-1}
\]
的第二项至多一次增长, 因而
\begin{equation}\label{eq:inverse-V-form}
V=4\ii\beta k^3\sg+k^2(-\ii\alpha\sg-4\beta Q)
+kW_1+W_0+\dfrac{\ii R}{k-a}.
\end{equation}
由 $\det\phi=\gamma$ 与 $x,t$ 无关, 两矩阵均迹零. 实谱上的共轭约化给出 $M^\dagger M=\gamma I$, 因此 $U,V$ 在实 $k\ne a$ 上反 Hermitian, 各系数具有相应约化.

作为同一 $\phi$ 的对数导数, 两矩阵恒满足零曲率关系. 比较 $k^2$ 项的非对角元及 $k$ 项的对角元, 得
\[
\begin{aligned}
(W_1)_{12}&=\alpha q-2\ii\beta q_x,\\
\partial_x(W_1)_{11}&=-2\ii\beta\partial_x|q|^2.
\end{aligned}
\]
故 $(W_1)_{11}=-2\ii\beta|q|^2+\chi_1(t)$. 由 \eqref{eq:direct-M-ends}, 两端 $V\to-FU_\infty$, 从而 $\chi_1(t)=0$. 再比较 $k$ 项的非对角元及常数项的对角元, 有
\[
\begin{aligned}
(W_0)_{12}&=2\beta|q|^2q+\dfrac{\ii\alpha}{2}q_x+\beta q_{xx},\\
\partial_x(W_0)_{11}&=\partial_x\left[\dfrac{\ii\alpha}{2}|q|^2-\beta(q\conj{q_x}-q_x\conj{q})\right].
\end{aligned}
\]
同一端态固定后一积分常数为 $\dfrac{\ii\Omega}{2}$. 于是 $W_1=V_1(q)$、$W_0=V_0(q)+\dfrac{\ii\Omega}{2}\sg$, 其余分量由迹零和约化确定, 得到完整规范 Lax 对 \eqref{eq:gaugelax}.

零曲率的共振留数和常数非对角元分别给出
\[
R_x=[-\ii a\sg+Q,R],\qquad
q_t=\ii\alpha\left(\dfrac{q_{xx}}{2}+|q|^2q\right)
+\beta(q_{xxx}+6|q|^2q_x)+2p+\ii\Omega q.
\]
前式为两条介质方程; 恢复 \eqref{eq:original-fields} 即得原 HMB 系统.

最后, 在普通连续谱上令
\[
\chi_-=Me^{-\ii\lambda x\sg},\qquad
\chi_+=M\diag(s_{22}^{-1},s_{22})e^{-\ii\lambda x\sg}.
\]
已恢复的空间方程及 \eqref{eq:direct-M-ends} 表明两者具有相应端点归一化, 由 Jost 唯一性即为重构场的 Jost 解. 所以 $S=\diag(s_{22}^{-1},s_{22})$, 反射为零, 透射函数为 \eqref{eq:onezeros}; 外部零点的列关系给出 $b_1(0)=c_1(0)s_{22}'(\zeta_1)=1$.
\end{proof}

\subsection{双曲表达与局域化}
写 $s=\xi+\ii\nu$, 光场公式中的投影元还可表示为
\begin{equation}\label{eq:hyperbolic-projector}
P_{12}=\dfrac{-2h_0\cosh(2\xi)+(1+h_0^2)\cos(2\nu)+\ii(1-h_0^2)\sin(2\nu)}
{2(1+h_0^2)\cosh(2\xi)-4h_0\cos(2\nu)}.
\end{equation}
分母不小于 $2(1-h_0)^2$, 局域包络与内部振荡分别由 $\xi$ 和 $\nu$ 控制.
